\section{讨论与结论}
\label{sec:conclusion}

FlashAttention-4针对非对称硬件扩展问题：由于张量核心的速度已非常快，主要瓶颈转移到了共享内存流量和指数运算吞吐量上，这促使我们进行算法和内核协同设计以缓解这些限制。我们围绕完全异步MMA重新设计了流水线，以将softmax与更大分块的矩阵乘法重叠，并引入了软件模拟指数运算和条件式softmax重缩放以减少非矩阵乘法操作。我们利用张量内存和2-CTA MMA模式减少共享内存流量。此外，2-CTA模式使全局原子累积的重构成为可能，将全局原子加操作的数量减半。FlashAttention-4完全使用嵌入Python的CuTe-DSL实现，在保持底层控制的同时实现了比基于C++模板的内核快20--30倍的编译速度。尽管针对Blackwell GPU进行了优化，随着计算能力持续超越非矩阵乘法单元，这些算法中的一些可以扩展到其他加速器。
